\documentclass{scrartcl}
\usepackage[a4paper, total={7in, 10in}]{geometry}
\usepackage{stmaryrd}
\usepackage{amsthm}
\usepackage{amssymb}
\usepackage{amsmath}
\usepackage{hyperref}
\usepackage{color}
\usepackage{tikz}
\usepackage{MacrosArt}

\title{Generalized Morse theory for distance functions}

\author{José Lorgeré}

\begin{document}

\maketitle

\begin{flushleft}

\tableofcontents

\newpage

\section{Generalities on Lipschitz Functions}

\subsection{Subdifferentials}

Let $f : U \subset \mathbb{R}^n \rightarrow \mathbb{R}^m$ be a Lipschitz function, where $U$ is
open. Rademacher's (source) theorem, which states that $f$ is differentiable almost everywhere, will
allow us to generalize differentials to this setting.

\begin{defi}
    The subdifferential of $f$ at $x$, denoted by $\partial f (x)$ is the convex hull of
    the linear maps $l$ obtained as
    \[ l = \lim_{i \rightarrow \infty} df(x_i) \]
    where $f$ is differentiable at each $x_i$ and $x_i \rightarrow x$.
\end{defi}

The subdifferential is then set valued. This set is always non empty and compact, $f$ being
Lipschitz bounds $df$ where it is defined.\newline
By (source Clarke), we can exclude a set of measure $0$ from the possible values of $x_i$
without changing the definition.

\begin{prop}
    $\partial f$ is the smallest convex valued upper semi-continuous function containing
    $df$ whenever it is defined.
\end{prop}

\begin{proof}
    First, $\partial f$ is upper semi-continuous : by compactness, is suffices to
    show that for all $x \in U, \varepsilon > 0$, for $y$ sufficiently close to $x$.
    \[ \partial f(y) \subset \partial f(x) + B(0, \varepsilon)\]
    Assume for contradiction that there exists a sequence $x_i \rightarrow x$ such that for all $i$,
    $\partial f (x_i)$ has an element out of $A := \partial f(x) + \overline{B}(0, \varepsilon/2)$.
    $A$ being convex and closed, let $x'_i$ be such that $df(x'_i)$ is out of
    $A$, and such that $x'_i - x_i \rightarrow 0$. Up to extraction, $df(x'_i)$ converges to an element in 
    $\partial f (x)$, so it enters $A$, contradiction.\newline \medskip
    The second part follow from $\partial f$ being a convex hull.
\end{proof}

We generalize the notion of Lipschitz functions to smooth manifolds

\begin{defi}
    A function $f : M \rightarrow \mathbb{R}^m$, where $M$ is a smooth manifold, is
    locally Lipschitz if for any chart $U \subset M$ and any compact $K \subset U$,
    there exists $C > 0$ such that, with the coordinates on $U$
    \[\forall x, y \in K, | f(x) - f(y) | \leq C |x - y| \]
\end{defi}

The following lemma generalizes the chain rule and also ensures that we can define
subdifferentials on manifolds unambiguously, i.e independently of the chart. For
$x \in M$, $\partial f(x)$ is then a compact convex subset of $T^*_x M$.

\begin{prop}
    Let $h : \mathbb{R}^m \rightarrow \mathbb{R}^k$ be a $C^1$ function. For $x \in U$
    \[ \partial (h \circ f)(x) \subset dh(f(x)) \circ \partial f(x) \]
    If $h$ is a diffeomorpism, we have equality.
\end{prop}

\begin{proof}
    For the first part, because $f$ is differentiable a.e, it suffices to show that
    the leftwise term is an upper semi-continuous convex valued function that contains
    $d(h \circ f)$ whenever $df$ is defined. This follows from the definition of $\partial f$.
    \medskip \newline
    If $h$ is a diffeomorphism, then for $x \in U$
    \[ \partial f(x) = \partial (h^{-1} \circ h \circ f)(x)
        \subset dh^{-1}(h \circ f)(x) \circ \partial (h \circ f)(x) \subset \partial f(x) \]
    so we have equalities, and the result follows.
\end{proof}

\subsection{The inverse function theorem for Lipschitz functions}

The subdifferentials allow us to generalize the idea of regular and critical points to
Lipschitz functions as follows

\begin{defi}
    A point $x \in U$ will be called critical if there exists some $l \in \partial f(x)$
    not of maximal rank. $x$ will be called regular if it isn't critical.
\end{defi}

We now consider the specific case of $f : U \subset \mathbb{R}^n \rightarrow \mathbb{R}^n$.
This section is devoted to proving a generalization by Clarke (source) of the inverse function
theorem for the case of Lipschitz functions.
\newline \newline
In this section, we norm the space of linear maps $\mathbb{R}^n \rightarrow \mathbb{R}^n$, which we denote
$\mathcal{L}(\mathbb{R}^n)$, with the operator norm $\| \dot  \|$ obtained from the norm used on
$\mathbb{R}^n$.

\begin{lem}\label{invfunlem}
    Assume $U$ is convex, and that there exists some closed convex set $C$ such that
    \begin{itemize}
        \item Every $l \in C$ is invertible.
        \item $\partial f(x) \subset C$ for all $x \in U$
    \end{itemize}
    Then there exists some $\delta > 0$ such that for all $x, y \in U$ :
    \[ |f(x) - f(y)| \geq \delta |x - y| \]
\end{lem}

\begin{proof}
    Firstly, we set
    \[ \delta = \inf\left\{ |l(h)| \mid l \in C, h \in S \right\} \]
    where $S$ is the unit sphere in $\mathbb{R}^n$. $\delta > 0$ because all $l \in C$
    are invertible and $C \times S$ is closed.
    \medskip \newline
    Let $x, y \in U$ and assume $x \neq y$. Let $\pi$ be the plane perpendicular to $y - x$
    passing through $x$. $f$ being differentiable a.e, Fubini's theorem implies that
    for almost all $x' \in \pi$, $f$ is differentiable almost on the whole line passing
    through $x'$ and perpendicular to $\pi$ (i.e $f$ isn't differentiable on a subset of this line of
    one dimensional measure $0$).\newline
    Denoting $F : t \longmapsto f(x' + t(y-x))$, we have :
    \begin{equation}\label{equ1}
        f(x' + y - x) - f(x') = \int_{0}^{1} F'(t) dt = \int_{0}^1 df(x' + t(y-x))(y-x)dt
    \end{equation}
    Define the linear map
    \[ L = \int_0^1 df(x'+t(y-x))dt \]
    $C$ being convex and closed, $L \in C$. Observe that the integral on the right of
    (\ref{equ1}) is just $L(y-x)$. So
    \[ |f(x'+y-x) - f(x')| \geq \delta |y-x| \]
    This inequality holds for almost all $x' \in \pi$, so by continuity of $f$, it holds for $x$
    and we get the desired result.
\end{proof}

\begin{theorem}
    If $x_0$ is a regular value of $f$, then there exists $V, W$, neighbourhoods of $x_0$ and $f(x_0)$
    respectively, such that $f$ induces a bijection $f : V \rightarrow W$ where $f^{-1}$ is
    Lipschitz too.
\end{theorem}

\begin{proof}
    By compactness of $\partial f(x_0)$, regularity of $x_0$, and the fact that the invertible
    linear maps form an open set, there exist some $\varepsilon > 0$ such that the set
    \[ C := \partial f(x_0) + \overline{B}(0, \varepsilon) \]
    is composed only of invertible maps. By upper semi-continuity of $\partial f$, there is
    some $r > 0$ such that for $y \in B(x_0, r)$, $\partial f(y) \subset C$. We then
    apply Lemma \ref{invfunlem} and get a corresponding $\delta > 0$.
    \medskip \newline
    Set $V = B(x_0, r/2)$ and $W = f(V)$. We showed in particular that $f$ is injective on
    $\overline{V}$, by compactness $f : \overline{V} \rightarrow f(\overline{V})$ induces a homeomorphism. $W$
    is then open. The inequality we got with $\delta$ implies that $f^{-1}$ is $\frac{1}{\delta}$
    Lipschitz.
\end{proof}

\end{flushleft}
 
\end{document}